\section{A specified arithmetic graph for the direct proof checker}
\label{sec:p05-pa-checker-arithmetization}

The direct tag-19 checker now receives a finite primitive recursive
description.  Its arithmetic formula is the graph of that description under
\cref{def:p05-pa-pr-graph}.  The construction also supplies an internal
\(\PA\) verification of modus ponens for this formula.  All proofs are in
\cref{annex:p5.checker-description,annex:p5.checker-internal-mp}.

\begin{hbdefinition}{def:p05-checker-description-language}{Finite description conventions}
The equations below are finite descriptions, with the displayed argument
order.  A reference to a previously defined function inserts its description;
application is composition, and
\[
 f(\bar x,0)=g(\bar x),\qquad
 f(\bar x,Sn)=h(\bar x,n,f(\bar x,n))
\]
is the primitive-recursion constructor with recursion argument last.  Auxiliary
parameters are listed in their order of first appearance.  Reordering or
duplicating arguments uses the corresponding projections.  A numeral constant
uses that many successors of zero.  Parentheses determine the composition
tree; displayed sums and products are left-associated.  A simultaneous state
is the cons-list of its fields, so every state update has one numerical output.

Use the already specified descriptions \(A,H,K,C\) of
\cref{ex:p05-pa-pair-cons-graphs}.  Add the following descriptions in order:
\[
\begin{aligned}
 D(0)&=0,&D(Sn)&=n,\\
 Z(0)&=1,&Z(Sn)&=0,\\
 M(x,0)&=0,&M(x,Sn)&=A(M(x,n),x),\\
 L(x,0)&=x,&L(x,Sn)&=D(L(x,n)),\\
 \mathsf{sg}(x)&=Z(Z(x)),&
 E(x,y)&=Z(A(L(x,y),L(y,x))),\\
 \mathsf{lt}(x,y)&=\mathsf{sg}(L(y,x)),&
 I(b,x,y)&=A(M(\mathsf{sg}(b),x),M(Z(b),y)).
\end{aligned}
\]
Thus \(D\) is predecessor and \(L\) is truncated subtraction.  In the
remaining specifications \(+\), multiplication, equality tests, strict-order
tests, negated tests, and Boolean conjunction/disjunction abbreviate
respectively \(A,M,E,\mathsf{lt},Z\), product, and sign of the sum.
Set \(\max(x,y)=I(\mathsf{lt}(x,y),y,x)\).
Set \(\min(x,y)=I(\mathsf{lt}(x,y),x,y)\); a maximum of more than two
arguments is left-associated.  A finite membership test is the sign of the
sum of its equality tests, in the displayed order.  A list of requirements
is the left-associated product of its tests in their printed order; an empty
list contributes one.  A bounded requirement at every position is the bounded
product in increasing position order.  Grammar alternatives are enumerated
in increasing constructor-tag order.  Repeated uses of a named expression
insert separate copies of its description, and local field names are expanded
by their indicated projections before composing the enclosing expression.
An ordered finite case distinction expands to nested \(I\), taking its first
applicable case, with its stated final default.

A constructor pattern \(c=N_k(b_0,\ldots,b_{r-1})\) first tests its tag
and arity.  A newly introduced field name abbreviates the corresponding
argument projection.  A field whose value was already specified contributes
an equality test between that projection and its specified expression.
Nested patterns are treated from outside inward and left to right.  Tuple
patterns use the same convention with a list-length test.  Constructing
\(N_k\) or a tuple as a value, rather than matching a pattern, uses its
displayed cons expression.  Thus pattern matching does not leave a choice
between different primitive recursive implementations.

A bounded fold with initial value \(s(\bar x)\), update
\(u(\bar x,j,v)\), and bound \(b(\bar x)\) means the description
\[
 V(\bar x,0)=s(\bar x),\quad
 V(\bar x,Sj)=u(\bar x,j,V(\bar x,j)),\quad
 V(\bar x,b(\bar x)).
\]
Bounds are evaluated before the fold; an update never changes its bound.
Bounded sums and products are this fold with initial values zero and one.
Least search over \(j<b\) starts with \([0,d]\), where \(d\) is the
specified failure value; at \(j\), replace \([0,v]\) by \([1,j]\)
exactly when the test succeeds, and otherwise retain the state.  Return its
second field.  All tests in a search are already defined descriptions.
Inclusive upper bounds mean successor bounds.  These clauses are expansion
instructions for each separately fixed finite expression, not a uniformly
total evaluator for arbitrary descriptions.
\end{hbdefinition}

\begin{hbdefinition}{def:p05-checker-list-description}{The fixed list and table library}
The notation \([a_0,\ldots,a_{r-1}]\) expands to
\(C(a_0,C(a_1,\ldots,C(a_{r-1},0)\ldots))\).  Define
\[
\begin{aligned}
 h(c)&=\sum_{a<c}a\,\mathsf{sg}\!\left(\sum_{b<c}E(C(a,b),c)\right),\\
 t(c)&=\sum_{b<c}b\,\mathsf{sg}\!\left(\sum_{a<c}E(C(a,b),c)\right),\\
 T(c,0)&=c,&T(c,Sn)&=t(T(c,n)),\\
 \ell(c)&=\sum_{j<c}\mathsf{sg}(T(c,j)),&
 e(c,j)&=h(T(c,j)),\\
 V(c,z,0)&=z,&V(c,z,Sn)&=C(e(c,n),V(c,z,n)),\\
 \mathsf{rev}(c)&=V(c,0,\ell(c)),&
 \mathsf{cat}(c,d)&=V(\mathsf{rev}(c),d,\ell(c)).
\end{aligned}
\]
Every natural number is a raw list code; \(h(0)=t(0)=0\), and entries
beyond a list's length have value zero.  Appending one entry means
\(\mathsf{cat}(c,[a])\).

For each fixed previously defined expression \(f(\bar x,a)\),
\(\mathsf{map}_f(\bar x,c)\) is the fold of length \(\ell(c)\)
starting at zero and appending \(f(\bar x,e(c,j))\) at step \(j\).
For a fixed zero-one test \(b(\bar x,a)\),
\(\mathsf{filter}_b\) uses the same fold, appending the unchanged entry
only when the test is one.  Set
\(\mathsf{range}(b)\) to the length-\(b\) fold starting at zero and
appending its counter.  Define
\[
 \mathsf{mem}(a,c)=\mathsf{sg}\!\left(\sum_{j<\ell(c)}E(a,e(c,j))\right),
 \qquad
 \mathsf{take}(c,n)=\mathsf{map}_{j\mapsto e(c,j)}
                   (\mathsf{range}(\min(n,\ell(c)))).
\]
The map notation in the last display takes \(c\) as an additional parameter.

A table with specified row expression \(u(\bar x,j,B)\) means
\[
 B(\bar x,0)=0,\qquad
 B(\bar x,Sj)=\mathsf{cat}
     (B(\bar x,j),[u(\bar x,j,B(\bar x,j))]).
\]
The row expression may consult only the given earlier-row list \(B\)
through \(e\); it does not call the table being defined.  Out-of-range
access consequently has the same zero default.  Write
\(k(c)=h(c)\), \(a_j(c)=e(t(c),j)\), and
\[
 \mathsf{node}_{k,r}(c)=\mathsf{sg}(c)E(h(c),k)E(\ell(t(c)),r),
 \qquad N_k(b_0,\ldots,b_{r-1})=C(k,[b_0,\ldots,b_{r-1}]).
\]
These are the actual constructor codes of the numerical checker.
\end{hbdefinition}

\begin{hbdefinition}{def:p05-checker-syntax-description}{Type and variable descriptions}
Define a table \(B(j)\) whose row at code \(c\) is
\([\mathsf{tm},\mathsf{fm},\mathsf{nn},\mathsf{mx}]\).
Use the constructor grammar of \cref{sec:arithmetic-coding-04}: tags
\(0,1,2,3,4\) are respectively a variable, zero, successor, addition,
multiplication; tags \(5,6,7,8,9,10,11,12\) are equality, bottom,
negation, conjunction, disjunction, implication, universal and existential
quantification.  Arity and child-type tests are exactly those in that grammar.
Child metadata are retrieved from the earlier-row list at the child code.
The term/formula fields are the zero-one disjunctions of their applicable
grammar cases.  Wrong tag, arity, or child type gives both fields zero.

At a valid node, \(\mathsf{nn}\) is one plus the sum of expression-child
node counts.  Its expression children, in order, are: none at tags \(0,1,6\);
\(a_0\) at tags \(2,7\); \(a_0,a_1\) at tags
\(3,4,5,8,9,10\); and \(a_1\) at tags \(11,12\).
The maximum-index field is \(a_0\) at a variable, zero at zero or bottom,
the maximum of child fields at a nonbinding compound, and
\(\max(a_0,\mathsf{mx}(a_1))\) at a quantifier.  Both metadata fields
are zero at an invalid node.  Thus the row expression uses only previously
defined list and numerical functions.  Retrieve row \(c\) from \(B(Sc)\)
to define \(\mathsf{Tm}(c),\mathsf{Fm}(c),n(c),m(c)\).

For each parameter \(i\), a second table has row
\([\mathsf{av},\mathsf{fv}]\) at code \(c\).  An invalid expression
has row \([0,0]\); a variable has both fields \(E(a_0(c),i)\); zero
and bottom have both fields zero.  Nonbinding compounds take the Boolean
disjunction of their children's corresponding fields.  At a quantifier with
binder \(b\) and matrix \(u\), its row is
\[
 [\ E(b,i)\lor\mathsf{av}(u),\quad
                  Z(E(b,i))\mathsf{fv}(u)\ ].
\]
The two children in this display mean entries from that table's earlier-row
list.  Retrieve row \(c\) after \(Sc\) updates to obtain
\(\mathsf{AV}(c,i),\mathsf{FV}(c,i)\).
Finally \(F(c)\) is the increasing-index fold through \(i<c\), appending
\(i\) exactly when \(\mathsf{Fm}(c)\mathsf{FV}(c,i)=1\), and
\[
 \mathsf{Sent}(c)\defeq\mathsf{Fm}(c)Z(\ell(F(c))).
\]
Every name here denotes its particular description and hence its particular
\(\Gamma\)-graph when used in arithmetic.
\end{hbdefinition}

\begin{hbdefinition}{def:p05-checker-substitution-description}{The fixed substitution and alpha machines}
The following two bounded folds fix the remaining syntax operations.  A state
is \([b,D,V]\): a success bit, a list of pending commands with the next
command first, and a value stack with its top first.  A failed state becomes
\([0,0,0]\) and stays there.  A successful state with \(D=0\) stays fixed.
Require exactly three state fields and \(b\in\{0,1\}\); any other state
fails before its commands are inspected.
Every command must have its specified length; every required value-stack entry
must exist.  Any other case fails.  Removing a command means taking \(t(D)\);
prepending a command sequence means concatenating that sequence to \(t(D)\).
Returning \(v\) means pushing \(v\) onto \(V\).

For substitution use commands
\[
 [0,c,i,u],\quad[1,c,j,z],\quad[2,k],\quad[3,k],
 \quad[4,k,j],\quad[5,k,z,i,u].
\]
They mean substitution, fresh free-variable renaming, unary assembly, binary
assembly, quantified assembly, and substitution after a completed renaming.
Their updates are the following ordered cases.
\begin{enumerate}
\item At \([0,c,i,u]\), require
 \((\mathsf{Tm}(c)\lor\mathsf{Fm}(c))\mathsf{Tm}(u)=1\).
 A variable returns \(u\) if its index equals \(i\), otherwise \(c\).
 Zero and bottom return \(c\).  A unary nonbinding node prepends its child's
 substitution command and \([2,k(c)]\).  A binary node prepends the left
 child's command, the right child's command, and \([3,k(c)]\), in that order.
 At \(Qx_j\theta\), return \(c\) if \(j=i\).  Otherwise, if
 \(\mathsf{AV}(u,j)=0\), prepend the matrix's substitution command and
 \([4,k(c),j]\).  In the remaining case let \(z\) be the least index
 at most \(1+\max(m(a_1(c)),m(u),i,j)\) absent from both all-variable
 tests and unequal to \(i,j\); use failure value zero.  Prepend
 \([1,\ulcorner\theta\urcorner,j,z]\) and \([5,k(c),z,i,u]\).
\item At \([1,c,j,z]\), require a valid expression.  A variable with
 index \(j\) returns \(N_0(z)\); another variable, zero, or bottom returns
 \(c\).  Nonbinding compounds recurse and assemble as in the first clause,
 using command tag 1 at their children.  A quantifier with binder \(j\)
 returns its entire code; another quantifier prepends its matrix's renaming
 command and the quantified assembly command with its unchanged binder.
 No further renaming occurs in this command mode.
\item Command \([2,k]\) is permitted for \(k\in\{2,7\}\); pop
 one value \(v\) and push \(N_k(v)\).  Command \([3,k]\) is permitted
 for \(k\in\{3,4,5,8,9,10\}\); pop \(v\), then \(w\), and push
 \(N_k(w,v)\).  Command \([4,k,j]\), for \(k\in\{11,12\}\),
 pops \(v\) and pushes \(N_k(j,v)\).
\item Command \([5,k,z,i,u]\), for \(k\in\{11,12\}\), pops the
 renamed matrix \(v\) and prepends \([0,v,i,u]\), then \([4,k,z]\).
\end{enumerate}
Start at \([1,[[0,c,i,u]],0]\), perform \(20(n(c)+1)^2\) updates,
and return the sole value exactly when the final state has success bit one,
empty command list, and a one-entry value stack; otherwise return zero.
Also return zero if the original expression/replacement type test fails.
This description is denoted \(\mathsf{sub}(c,i,u)\).

For alpha comparison use the same state format, with commands
\([0,c,E],[1,k],[2,k],[3,k]\).  The first command traverses \(c\)
under the list \(E\) of enclosing binder indices, nearest first.  At a
variable with index \(i\), search \(j<\ell(E)\) for the first
\(e(E,j)=i\), with default \(\ell(E)\).  Return \(N_{18}(j)\)
when found and \(N_{17}(i)\) otherwise.  Zero and bottom return their codes.
A unary or binary nonbinding node prepends its child commands in order and
assembly command 1 or 2; these assemble \(N_k(v)\) or \(N_k(w,v)\),
with the same pop order and permitted tags as above.  A quantifier with binder
\(i\) prepends its matrix command under \(C(i,E)\) and \([3,k]\);
this last command, for \(k\in\{11,12\}\), pops \(v\) and pushes
\(N_k(v)\), with one argument.  Invalid expressions or commands fail.
Start at \([1,[[0,c,0]],0]\), use \(4(n(c)+1)\) updates, and extract
the final singleton by the same rule.  Require \(\mathsf{Fm}(c)=1\)
at entry; otherwise return zero.  Call the result \(\mathsf{ar}(c)\), and set
\[
 \mathsf{ae}(c,d)\defeq
          \mathsf{Fm}(c)\mathsf{Fm}(d)E(\mathsf{ar}(c),\mathsf{ar}(d)).
\]
All command decoding, stack changes, bounds, and searches expand through the
preceding library.  These are finite machine transition descriptions with a
fixed number of updates, not instructions to run an unbounded computation.
\end{hbdefinition}

\begin{hbdefinition}{def:p05-checker-axiom-description}{The direct axiom description}
Define \(U(P,c)\) by folding over \(\mathsf{rev}(P)\), starting at
\(c\) and replacing the state \(v\) at index \(j\) by
\(N_{11}(e(\mathsf{rev}(P),j),v)\).  Use exactly the constructor expression
for \(J(m,i)\) and \(\mathsf{IndAx}(m,i)\) in
\cref{def:p05-canonical-induction-code}, with the descriptions \(F,U,\mathsf{sub}\)
just specified, and with the filter retaining indices unequal to \(i\).
Retain the zero default when \(\mathsf{Fm}(m)=0\).

The prefix-stripping fold starts at \([a,0]\) and has bound \(a\).
On state \([c,P]\), if \(\mathsf{node}_{11,2}(c)=1\), replace it by
\([a_1(c),\mathsf{cat}(P,[a_0(c)])]\); otherwise retain it.
From its final \([c,P]\), recover the fields of
\[
 c=N_{10}(d,N_{11}(i,m)),\qquad
 d=N_8(b,N_{11}(i',N_{10}(m',s))).
\]
Every displayed tag and arity is tested by \(\mathsf{node}\).
Let \(B(a)\) be the conjunction of these tests,
\(E(i,i'),E(m,m'),\mathsf{Fm}(m)\), and
\(E(a,\mathsf{IndAx}(m,i))\).  Field extraction is total even if a test
fails.  Define
\[
\begin{aligned}
 Q_7&\defeq[q_1,q_2,q_3,q_4,q_5,q_6,q_7],\\
 \mathsf{Ax}_{19}(a)&\defeq\mathsf{Sent}(a)
       \mathsf{sg}\!\left(B(a)+\sum_{j<7}E(a,e(Q_7,j))\right).
\end{aligned}
\]
The constant \(Q_7\) expands to the displayed seven-entry cons-list, and
\(e(Q_7,j)\) uses the previously fixed total entry description, including
its zero value when \(j\geq7\).  The sum uses the declared bounded fold
with this particular selector.  Each of the seven \(q_j\) constants inserts
its exact finite constructor expression in
\cref{sec:p05-standard-pa-checker}.  Thus no axiom code or matrix remains an
unspecified constant or a search over formulas.
\end{hbdefinition}

\begin{hbdefinition}{def:p05-checker-row-description}{The exact row description}
For a raw row list \(R\), put \(r_j=e(R,j)\), and use its five fields
\([r,a,J,U,l]\) by entry projection.  Write \(f_R(j)=e(r_j,1)\) and
\(\lambda_R(j)=e(r_j,4)\).  Form a table of pairs \([b_j,O_j]\),
starting empty and appending one pair at each \(j<\ell(R)\).
For a premise position \(v\), abbreviate
\[
 j_v=e(J,v),\qquad f_v=f_R(j_v),\qquad
 O_v=e(e(B,j_v),1),
\]
where \(B\) is the earlier-row table.  Missing entries have value zero.
For parameters \(i,\theta\), define
\[
 D_R^{i,\theta}(O)\defeq
 \mathsf{filter}_{z\mapsto
    Z(E(\lambda_R(z),Si)E(f_R(z),\theta))}(O).
\]
Its nonvacuity test is \(\mathsf{lt}(\ell(D_R^{i,\theta}(O)),\ell(O))\).

The common row test requires length five, \(r<20\),
\(\mathsf{Fm}(a)=1\), the following premise and parameter arities, and
\(j_v<j\) at every premise position:
\[
\begin{aligned}
 &(0,2,1,1,1,1,3,1,2,1,2,1,1,1,1,1,2,1,0,2),\\
 &(0,0,0,0,0,0,2,1,0,1,0,0,1,3,3,1,4,0,1,4).
\end{aligned}
\]
For \(r\ne0\) it also requires \(l=0\).  Parentheses in these two
displays specify fixed lookup lists under the cons coding.

Here is the complete additional test and open-list update for each tag.
Formula letters are recovered by constructor projections; every constructor
pattern includes its exact tag and arity.  Parameters displayed as terms or
formulas receive \(\mathsf{Tm}\) or \(\mathsf{Fm}\) tests.  All formula
equalities use \(E\).  Set \(O=0\) in a case with a wrong tag or arity.
\begin{enumerate}
\item[0.] No additional test; \(O=[j]\).
\item[1.] \(a=N_8(f_0,f_1)\); \(O=\mathsf{cat}(O_0,O_1)\).
\item[2,3.] \(f_0=N_8(c,d)\); \(a=c\) for tag 2 and \(a=d\)
 for tag 3; \(O=O_0\).
\item[4,5.] \(a=N_9(c,d)\); \(f_0=c\) for tag 4 and \(f_0=d\)
 for tag 5; \(O=O_0\).
\item[6.] \(U=[i,k]\), \(f_0=N_9(c,d)\), \(f_1=f_2=a\).
 Both \(D_R^{i,c}(O_1)\) and \(D_R^{k,d}(O_2)\) must be nonvacuous;
 \(O=\mathsf{cat}(O_0,\mathsf{cat}(D_R^{i,c}(O_1),D_R^{k,d}(O_2)))\).
\item[7.] \(U=[i]\), \(a=N_{10}(c,f_0)\); require nonvacuity and put
 \(O=D_R^{i,c}(O_0)\).
\item[8.] \(f_1=N_{10}(f_0,a)\);
 \(O=\mathsf{cat}(O_0,O_1)\).
\item[9.] \(U=[i]\), \(a=N_7(c)\), \(f_0=N_6()\);
 require nonvacuity and put \(O=D_R^{i,c}(O_0)\).
\item[10.] \(f_1=N_7(f_0)\), \(a=N_6()\);
 \(O=\mathsf{cat}(O_0,O_1)\).
\item[11.] \(f_0=N_6()\); \(O=O_0\).
\item[12.] \(U=[i]\), \(f_0=N_6()\); require nonvacuity and put
 \(O=D_R^{i,N_7(a)}(O_0)\).
\item[13.] \(U=[i,c,u]\) with formula \(c\) and term \(u\);
 \(f_0=N_{11}(i,c)\), \(a=\mathsf{sub}(c,i,u)\); \(O=O_0\).
\item[14.] The same parameter types; \(f_0=\mathsf{sub}(c,i,u)\),
 \(a=N_{12}(i,c)\); \(O=O_0\).
\item[15.] \(U=[i]\), \(a=N_{11}(i,f_0)\), and
 \(\prod_{v<\ell(O_0)}Z(\mathsf{FV}(f_R(e(O_0,v)),i))=1\);
 \(O=O_0\).
\item[16.] \(U=[i,c,y,k]\) with formula \(c\),
 \(f_0=N_{12}(i,c)\), \(f_1=a\).
 Put \(O'=D_R^{k,\mathsf{sub}(c,i,N_0(y))}(O_1)\), require its
 nonvacuity, and put \(O=\mathsf{cat}(O',O_0)\).  Require
 \[
 Z(\mathsf{AV}(a,y))Z(\mathsf{AV}(f_0,y))
       \prod_{v<\ell(O)}Z(\mathsf{AV}(f_R(e(O,v)),y))=1.
 \]
\item[17.] \(\mathsf{ae}(f_0,a)=1\); \(O=O_0\).
\item[18.] \(U=[u]\) with term \(u\), \(a=N_5(u,u)\); \(O=0\).
\item[19.] \(U=[u,v,c,i]\) with terms \(u,v\) and formula \(c\);
 \(f_0=N_5(u,v)\), \(f_1=\mathsf{sub}(c,i,u)\),
 \(a=\mathsf{sub}(c,i,v)\); \(O=\mathsf{cat}(O_0,O_1)\).
\end{enumerate}
The row's bit is the conjunction of its common and additional tests.  If
this bit is zero, store \([0,0]\); otherwise store \([1,O]\).
The preceding cases, finite lists, folds, and projections therefore determine
one finite row-table description \(\mathsf{Rows}(R)\).

Define the parent count and record test by
\[
\begin{aligned}
 \mathsf{pc}(R,i)&=\sum_{j<\ell(R)}
                    \sum_{v<\ell(e(r_j,2))}E(e(e(r_j,2),v),i),\\
 \mathsf{Tree}(R)&=\mathsf{sg}(\ell(R))
       \prod_{j<\ell(R)}e(e(\mathsf{Rows}(R),j),0)
       \prod_{i<\ell(R)}E(\mathsf{pc}(R,i),Z(E(i,D(\ell(R))))).
\end{aligned}
\]
For \(\ell(R)>0\), put
\(\mathsf{root}(R)=f_R(D(\ell(R)))\) and
\(\mathsf{open}(R)=e(e(\mathsf{Rows}(R),D(\ell(R))),1)\);
both have default zero for the empty record.
\end{hbdefinition}

\begin{hbdefinition}{def:p05-selected-direct-pa-graph}{The selected direct arithmetic proof formula}
Decode \(p\) by requiring \(p=N_{19}(d,0)\), \(d=N_{15}(R)\),
with exact arities.  The binary characteristic description \(\chi_{19}(p,x)\)
returns zero if these outer tests fail.  Otherwise it returns
\[
 \mathsf{Tree}(R)\mathsf{Sent}(\mathsf{root}(R))E(x,\mathsf{root}(R))
     \prod_{v<\ell(\mathsf{open}(R))}
       \mathsf{Ax}_{19}(f_R(e(\mathsf{open}(R),v))).
\]
Let \(d_{19}\) be the finite formation tree obtained by the stated
expansions of this description and all its preceding library calls.  Select
the actual arithmetic formulas
\[
 \mathsf{Prf}^{19}_{\PA}(p,x)\defeq\Gamma_{d_{19}}(p,x,\bar1),
 \qquad
 \mathsf{Prv}^{19}_{\PA}(x)\defeq
                   \exists p\,\mathsf{Prf}^{19}_{\PA}(p,x).
\]
The \(\Gamma\) construction uses its fixed fresh-variable convention at
every expansion.  These formulas are selected by a finite description,
without selecting an arbitrary formula having the same standard values.
\end{hbdefinition}

\begin{hbtheorem}{thm:p05-selected-direct-pa-graph}{proved=annex:p5.checker-description}{Graph meaning and internal library equations}
The preceding specifications determine finite primitive recursive descriptions.
Numerically \(\chi_{19}\) is the characteristic function of
\(\mathsf{StdProof}_{\PA}\) from
\cref{def:p05-standard-pa-numerical-proof}.  The selected formulas are
\(\Sigma^0_1\), and for every arithmetic sentence \(\sigma\),
\[
 \PA\derives\sigma
 \quad\Longleftrightarrow\quad
 \mathcal N\entails
       \mathsf{Prv}^{19}_{\PA}(\overline{\ulcorner\sigma\urcorner}).
\]
For their particular graph descriptions, \(\PA\) proves the list,
table, and syntax equations used in
\cref{annex:p5.checker-description}, including constructor decoding,
concatenation entries, and
\[
 \mathsf{Sent}(N_{10}(u,v))=1
       \quad\longrightarrow\quad
            \mathsf{Sent}(u)=\mathsf{Sent}(v)=1.
\]
All function notation in this last assertion abbreviates existentially
quantified graph outputs, with uniqueness supplied by
\cref{thm:p05-pa-pr-graphs}.
\end{hbtheorem}

\begin{hbdefinition}{def:p05-direct-mp-description}{The fixed direct modus-ponens description}
For an offset \(a\), let \(\mathsf{sh}(a,J)\) map each entry \(j\)
of \(J\) to \(a+j\).  Let \(\mathsf{sr}(a,r)\) return
\([b,c,\mathsf{sh}(a,J),U,l]\) when \(r=[b,c,J,U,l]\) has length
five, and return zero otherwise.  Let \(\mathsf{SR}(a,R)\) map
\(\mathsf{sr}(a,\cdot)\) over \(R\).

To define \(\mathsf{mp}_{19}(p,q)\), decode both tag-19 records as above,
write their row lists as \(R_p,R_q\), and their roots as \(c_p,c_q\).
Require \(\chi_{19}(p,c_p)=\chi_{19}(q,c_q)=1\) and
\(c_p=N_{10}(c_q,z)\) with exact arity.  Failure returns zero.
On success set \(a=\ell(R_q)\), \(b=\ell(R_p)\), and return
\[
 N_{19}\!\left(N_{15}\!\left(
   \mathsf{cat}\!\left(R_q,
     \mathsf{cat}(\mathsf{SR}(a,R_p),
          [[8,z,[D(a),D(a+b)],0,0]])\right)\right),0\right).
\]
This inserts the antecedent proof first and the implication proof second,
as required by the actual implication-elimination rule.  Call its expanded
description \(d_{\rm mp}\).  It includes the displayed failure test.
\end{hbdefinition}

\begin{hbtheorem}{thm:p05-direct-pa-internal-mp}{proved=annex:p5.checker-internal-mp}{Internal modus ponens and D2 for the selected formula}
In \(\PA\), with all free numerical variables universally quantified,
\[
\begin{split}
 &\mathsf{Prf}^{19}_{\PA}(p,N_{10}(u,v))
       \land\mathsf{Prf}^{19}_{\PA}(q,u)\\
 &\quad\longrightarrow\exists r\,
       (\Gamma_{d_{\rm mp}}(p,q,r)
                         \land\mathsf{Prf}^{19}_{\PA}(r,v)).
\end{split}
\]
The constructor in the first line is eliminated through its specified graph.
Consequently, for every fixed pair of arithmetic sentences \(\varphi,\psi\),
\[
\begin{split}
 \PA\derives\;&\mathsf{Prv}^{19}_{\PA}
                   (\overline{\ulcorner\varphi\to\psi\urcorner})\\
 &\to\bigl(\mathsf{Prv}^{19}_{\PA}
                    (\overline{\ulcorner\varphi\urcorner})
       \to\mathsf{Prv}^{19}_{\PA}
                    (\overline{\ulcorner\psi\urcorner})\bigr).
\end{split}
\]
Also, if \(\PA\derives\sigma\), then
\(\PA\derives\mathsf{Prv}^{19}_{\PA}
                     (\overline{\ulcorner\sigma\urcorner})\).
This last metatheoretic implication is D1 for this selected \(\PA\) formula.
\end{hbtheorem}

\begin{hbexample}{def:p05-direct-graph-empty-proof}{The empty-context identity proof in the selected graph}
Let \(o=N_1()\), \(e=N_5(o,o)\), and
\(p=N_{19}(N_{15}([[18,e,0,[o],0]]),0)\).
The specified description gives \(\chi_{19}(p,e)=1\), so its finite
numeral computation yields
\(\PA\derives\mathsf{Prf}^{19}_{\PA}(\bar p,\bar e)\).
The direct-axiom product is empty in this computation.  No assertion that
\(0=0\) is a nonlogical axiom is used.
\end{hbexample}

\paragraph{Evidence-to-proof construction and positive introspection.}
For each fixed closed \(\Sigma^0_1\) sentence \(\sigma\) in its
original syntax, \cref{thm:p05-original-sigma-evidence} supplies the
specified unary description \(\widehat d_\sigma^+\) and proves in
\(\PA\)
\[
 \forall e\,\bigl(\mathsf{Ev}_\sigma(e)\to
    \exists p\,(\Gamma_{\widehat d_\sigma^+}(e,p)\land
       \mathsf{Prf}^{19}_{\PA}(p,\overline{\ulcorner\sigma\urcorner}))\bigr).
\]
The evidence is exactly the formula constructed in
\cref{thm:p05-pa-sigma-evidence}; the signed proof constructors retain
the original sentence code rather than ending at its normal form.
\Cref{thm:p05-literal-graph-evidence,thm:p05-direct-pa-d3} applies
the construction to the literal selected provability sentence and proves D3.
Together with the internal modus-ponens theorem, this supplies D1--D3
for the same selected graph. An earlier arbitrarily chosen numeralwise
representative is a separate arithmetic formula; no internal identification
with it is assumed. The full later provability-logic completeness
constructions retain their own proof obligations.
